\section{Related Work}
\subsection{Epidemiology and prior KNHANES work}
Tinnitus and hearing loss often coexist but relate heterogeneously. KNHANES
analyses have estimated Korean tinnitus prevalence and correlates
\citep{park2014tinnitus,joo2015hrqol}, and NHANES has enabled U.S. studies of
audiometry, occupational noise, tinnitus, and mental-health correlates
\citep{hoffman2017declining,chakrabarty2024depression,mahboubi2013noise}; reviews
emphasize heterogeneous case definitions and the need for harmonized variables
when pooling across surveys \citep{mccormack2016review}. Evidence linking
psychological stress to tinnitus is consistent but largely observational, with
residual confounding by noise, sleep, and cardiometabolic factors difficult to
exclude \citep{baigi2011tinnitus}. On the same KNHANES cycles, prior work has
examined tinnitus alongside distress---most closely \citet{lee2018tinnitus}, on
higher perceived stress and mood symptoms in premenopausal women with tinnitus.
STARS is therefore \emph{not} the first to link tinnitus with distress in KNHANES;
its contribution is methodological: (i) a \emph{prespecified, falsifiable} contrast
between the tinnitus \emph{symptom} and the audiometric \emph{threshold} (H1)
rather than a single tinnitus--distress association; (ii) an explicit DAG that
separates confounders from mediators, reporting both minimal-sufficient and
mediator-adjusted models; (iii) a cross-national directional-consistency check in
NHANES under a related distress construct; and (iv) a fully reproducible open
scaffold. The novelty is design discipline and reproducibility, not the mere
existence of a tinnitus--stress association.

\subsection{Beyond the pure-tone average}
Because real-world listening difficulty and tinnitus burden can persist when the
better-ear pure-tone average appears near-normal, STARS separates tinnitus
presence from bothersome tinnitus and analyzes worse-ear and high-frequency
thresholds and asymmetry, not better-ear PTA alone.

\subsection{Clinical pathways and reporting standards}
Practice guidelines frame the endpoints. The tinnitus guideline distinguishes
bothersome from non-bothersome tinnitus and prioritizes education,
cognitive-behavioral strategies, and sound therapy \citep{tunkel2014cpg}, while the
sudden-hearing-loss guideline defines a ``golden window''---prompt audiometric
confirmation and early corticosteroid therapy within roughly two weeks---so delayed
presentation is associated with worse outcomes \citep{chandrasekhar2019ssnhl}. This
is the clinical rationale for the deterministic red-flag referral-safety layer
developed in the companion SAFE-EAR paper. The observational analyses adhere to
STROBE \citep{vonelm2007strobe}; the prospective prediction and open-medical-LLM
components follow the TRIPOD+AI statement \citep{collins2024tripodai}.
